\documentclass{article}
\usepackage{propositions}

%% The built-in `fitmargin' style, and the environment-level machinery it rests
%% on: an item-level key set for the environment, and \propformlabel.
%%
%% The rules mark the surrounding text margins.  In every list below the label
%% must start at the left rule and the body must start one \labelsep past the
%% end of the label, with continuation lines aligned under the body.

\newcommand\dims{\texttt{LM=\the\leftmargin\ LW=\the\labelwidth\
  LS=\the\labelsep}}

\begin{document}

\section{The style itself}

\noindent\rule{\textwidth}{0.3pt}
\begin{prop}[name=Name of the item, style=fitmargin]
  \item Text of the item, which is long enough that we can see what happens
  when it wraps to the next line.
\end{prop}
\noindent\rule{\textwidth}{0.3pt}
\begin{prop}[name=A considerably longer name for the item, style=fitmargin]
  \item Text of the item, which is long enough that we can see what happens
  when it wraps to the next line.
\end{prop}
\noindent\rule{\textwidth}{0.3pt}

The margin is sized to the label, so the second list is indented further than
the first.  \texttt{LW} must equal the width of the label and \texttt{LM} must
be \texttt{LW+LS}; if \texttt{LW} kept the class default (20pt) the label would
overflow its box and start short of the margin.

\begin{prop}[name=Short, style=fitmargin]
  \item \dims
\end{prop}

\section{Numbered items and \texttt{items}}

Without \texttt{items} the width is computed from the first label; with it,
from the \textit{n}th.  The second list below must therefore be the wider one,
and both must number from 1.

\begin{prop}[style=fitmargin]
  \item First item, long enough to wrap onto a second line for comparison.
  \item Second item. \label{h:a}
\end{prop}

\begin{prop}[style=fitmargin, items=20]
  \item First item, long enough to wrap onto a second line for comparison.
  \item Second item. \label{h:b}
\end{prop}

References: \ref{h:a}, \ref{h:b}.  Expected: (2), (4) --- unnamed items keep
the format of the \texttt{numbered} style, parentheses and all; the style only
sizes the margin.  (Numbered items are not what this style is for; they are
used here to exercise \texttt{items} and the preview.)

\section{The preview must not disturb the counter}

Computing \verb|\propformlabel| resolves an item, which steps the counter, and
the step has to be undone --- but only when it happened.  A named environment
does not step at all, so an unconditional undo would run the counter backwards:
one step per environment.

\begin{prop}[name=Named, style=fitmargin]
  \item This environment's items are named, so nothing is stepped.
  \item Nor here.
\end{prop}

\begin{prop}[style=fitmargin, items=20]
  \item Numbered, after a named environment. \label{h:c}
  \item Second. \label{h:d}
\end{prop}

\begin{prop}
  \item Numbered, plain environment. \label{h:e}
\end{prop}

References: \ref{h:c}, \ref{h:d}, \ref{h:e}.  Expected: (5), (6), (7) ---
consecutive, with nothing lost to the preview.

\section{The margin does not jitter}

\verb|fitmargin| measures \verb|\propwidestlabel|, not \verb|\propformlabel|, so
the margin depends on the number of digits in the label and not on which digits
they are.  In the four lists below \texttt{LM} must be \emph{equal} in the first
pair (both labels two digits) and strictly increasing from the second pair to
the third and fourth (two, then three, then four digits).  Were the actual
label measured, the first pair would differ.

\setcounter{numpropi}{10}
\begin{prop}[style=fitmargin]
  \item Label (11): \dims
\end{prop}
\setcounter{numpropi}{98}
\begin{prop}[style=fitmargin]
  \item Label (99): \dims
\end{prop}
\setcounter{numpropi}{998}
\begin{prop}[style=fitmargin]
  \item Label (999): \dims
\end{prop}
\setcounter{numpropi}{9998}
\begin{prop}[style=fitmargin]
  \item Label (9999): \dims
\end{prop}
\setcounter{numpropi}{0}

A label with no digits is unaffected, and formatting survives the widening:
\verb|\propwidest| of \textbf{(3.147)} is \propwidest{\textbf{(3.147)}}, which
must still be bold, and of \textsc{Prop}~1a is \propwidest{\textsc{Prop}~1a},
which must still be small caps with the space intact.

\section{Labels that change width within one list}

\verb|\propformlabel| predicts the label this environment's items will carry,
and keeps that value for the whole environment.  It must not be updated to the
label just produced: a style's dimension keys are read again for every item's
label area, so a changing prediction sizes each item's label box to its
\emph{predecessor's} label, and a label that has grown a digit then overflows
and shoves its own first line sideways.

Nothing else here catches that.  A named environment gives every item the same
label, and \verb|items=20| above numbers two items \texttt{(2)} and
\texttt{(4)} --- equal widths either way.  It takes a list that crosses a
digit boundary.

\setcounter{numpropi}{8}
\noindent\rule{\textwidth}{0.3pt}
\begin{prop}[items=2, style=fitmargin]
  \item Label (9), one digit.
  \item Label (10), two digits: this item's body must start at exactly the
  same place as the one above.
\end{prop}
\noindent\rule{\textwidth}{0.3pt}
\setcounter{numpropi}{0}

Both bodies must begin at the same horizontal position, and both labels at the
left rule.  With the prediction refreshed per item, the second body sat
4.6pt --- the width difference between \texttt{(9)} and \texttt{(10)} ---
further right than the first.

\section{\texttt{outerfit}: the outer level only}

\verb|outerfit| fits the margin at level one and does nothing below it, so that
it can be set for a whole document with \verb|\propoptions{style=outerfit}|
without disturbing nested lists.  Its dimension values vary with the level via
\verb|\proplevelchoice|, the second entry being the sentinel \verb|*|.  It
differs from \verb|fitmargin| in separating label from text by 1.3em rather
than a \verb|\labelsep|, \textsf{linguex}'s spacing for examples.

{\propoptions{style = outerfit}
\noindent\rule{\textwidth}{0.3pt}
\begin{prop}[name=Outer Name]
  \item Outer, named: \dims
  \begin{prop}[style=roman]
    \item Nested, roman: \dims
    \begin{prop}
      \item Third level: \dims
    \end{prop}
  \end{prop}
\end{prop}
\begin{prop}
  \item Outer, numbered: \dims
\end{prop}
\noindent\rule{\textwidth}{0.3pt}
\par}

The two outer items must be fitted to their own labels: \texttt{LM} equal to
the label's own width plus 1.3em (13pt at 10pt), and the label starting at the
left rule.  The label \emph{box} is wider than the label by the difference
between 1.3em and \texttt{LS}, which is invisible while the label is
left-aligned in it.  The nested ones must keep the class
defaults untouched --- \texttt{LM=22.0pt} at level two
(\verb|\leftmarginii|) and \texttt{LM=18.7pt} at level three
(\verb|\leftmarginiii|) in \texttt{article} --- which is the whole point of the
style, and what a plain \verb|style=fitmargin| set globally would not do.

Since the setting is confined to a group above, the lists below are unaffected.

A level-two list is not always at \LaTeX{} list depth two: a box resets the
depth, and at depth one the class sets no \verb|\labelwidth| at all.  The
sentinel must still leave a sane value there, not whatever the enclosing list
left behind --- which under \verb|outerfit| would be the outer item's fitted
width.

{\propoptions{style = outerfit}
\noindent\rule{\textwidth}{0.3pt}
\begin{prop}
  \item Outer: \dims
  \par\noindent
  \begin{minipage}{0.7\textwidth}
    \begin{prop}
      \item Depth one: \dims
    \end{prop}
  \end{minipage}
\end{prop}
\noindent\rule{\textwidth}{0.3pt}
\par}

The inner list must show \texttt{LM=25.0pt} and \texttt{LW=20.0pt}, the depth-one
defaults, exactly as it would without \verb|outerfit| --- not the outer item's
\texttt{LW}.

\section{An environment-level name reaches the items}

\texttt{fitmargin} is only useful because a name set for the environment applies
to each of its items.  A blank \verb|\item| argument must leave that name alone
rather than override it with an empty one.

\begin{prop}[name=Inherited]
  \item No argument at all. \label{h:f}
  \item[] An empty argument. \label{h:g}
  \item[gloss=g] An argument with no name in it. \label{h:h}
  \item[Own name] Its own name, which wins. \label{h:i}
\end{prop}

References: \ref{h:f}, \ref{h:g}, \ref{h:h}, \ref{h:i}.

Expected references: \textbf{Inherited}, \textbf{Inherited},
\textbf{Inherited}, \textbf{Own name} --- all in the named style, since a name
reaches every one of them, three from the environment and one from the item's
own argument.  The third item's \emph{label} carries the gloss as well,
\textbf{Inherited}~(g), but a gloss is not part of the reference.

\end{document}
